%% kapa: exam=1 round=34 year=2023 period=2 form=A subject=appraisal_law subject_ko=감정평가관계법규 seq=11 number=11
%% source: Q-Net 감정평가사 기출문제 (cst00309) artlSeq=5211789 file=raw/1008/5211789/2023년도 제34회 감정평가사 1차시험 2교시.pdf
%% conversion: kordoc 4.24.51 -> tools/gen1.py (auto); review=proofread-sonnet; flags=none
\begin{question}{11}
국토의 계획 및 이용에 관한 법령상 도시·군계획시설사업의 단계별 집행계획에 관한 설명으로 옳은 것은?
\begin{choices}
  \item 국토교통부장관은 단계별 집행계획의 수립주체가 될 수 있다.
  \item 「도시 및 주거환경정비법」에 따라 도시·군관리계획의 결정이 의제되는 경우에는 해당 도시·군계획시설결정의 고시일부터 3년 이내에 단계별 집행계획을 수립할 수 있다.
  \item 단계별 집행계획은 제1단계 집행계획과 제2단계 집행계획 및 제3단계 집행계획으로 구분하여 수립한다.
  \item 3년 이내에 시행하는 도시·군계획시설사업은 제2단계 집행계획에 포함되도록 하여야 한다.
  \item 단계별 집행계획이 수립되어 공고되면 변경할 수 없다.
\end{choices}
\end{question}
\answer{1}
